\documentclass[10pt,a4paper]{extarticle}
\usepackage[utf8]{inputenc}
\usepackage[T1]{fontenc}
\usepackage[spanish,es-noshorthands,es-nodecimaldot]{babel}
\usepackage{lmodern}
\usepackage[margin=0.65cm]{geometry}
\usepackage{amsmath,amssymb,multicol,enumitem,xcolor,array,booktabs}
\usepackage[most]{tcolorbox}
\definecolor{azul}{HTML}{1F4E79}
\definecolor{azulclaro}{HTML}{E8F0F8}
\definecolor{naranja}{HTML}{B85C00}
\pagestyle{empty}
\setlength{\parindent}{0pt}\setlength{\parskip}{0.8pt}
\setlist{nosep,leftmargin=1em}
\setlength{\columnsep}{0.32cm}\setlength{\columnseprule}{0.2pt}
\renewcommand{\arraystretch}{1.05}
\newcommand{\T}[1]{\par\vspace{2.5pt}{\color{white}\colorbox{azul}{\parbox{\dimexpr\linewidth-2\fboxsep}{\bfseries\MakeUppercase{#1}}}}\par\vspace{1pt}}
\newcommand{\st}[1]{{\color{azul}\bfseries #1}}
\newcommand{\ojo}[1]{{\color{naranja}\bfseries #1}}
\newcommand{\imp}{\Rightarrow}\newcommand{\sii}{\Leftrightarrow}
\newcommand{\p}[1]{\ensuremath{\textit{#1}}}
\newtcolorbox{caja}{colback=azulclaro,colframe=azul,boxrule=0.4pt,arc=1pt,left=2pt,right=2pt,top=1pt,bottom=1pt}
\begin{document}
% =============================== FRENTE ===============================
\begin{multicols*}{3}
{\large\bfseries\color{azul} IA · Examen 1 · Formulario (frente)}\par

\T{Agentes y modelado del ambiente}
\st{Agente racional}: elige la acción que maximiza el desempeño \emph{esperado} dada la secuencia de
percepciones y su conocimiento (no es omnisciente: lo esperado, no lo real).
\st{PEAS}: \emph{P}erformance (desempeño), \emph{E}nvironment, \emph{A}ctuadores, \emph{S}ensores.
\st{Modelar}: estados, perceptos, acciones, transición, desempeño.
\st{Propiedades} (justificar cada una en 1 línea):
\begin{itemize}
\item observable total / \emph{parcial} (¿el sensor ve todo el estado relevante?)
\item un agente / \emph{multiagente} (cooperativo, competitivo)
\item determinista / \emph{estocástico} (¿el resultado de la acción es seguro?)
\item episódico / \emph{secuencial} (¿la decisión afecta el futuro?)
\item estático / \emph{dinámico} (¿cambia mientras el agente piensa?)
\item discreto / \emph{continuo}; conocido / desconocido
\end{itemize}
Mundo real (taxis, bicis, robots): parcial, multi, estocástico, secuencial, dinámico, continuo (discretizar).
\st{Tipos}: reflejo (regla condición–acción, sin consecuencias), con modelo, basado en metas
(«¿qué pasa si…?»), basado en utilidad. Cuatro enfoques: pensar/actuar × humano/racional; el curso:
\emph{actuar racionalmente}. Turing (conducta), cuarto chino (sintaxis $\ne$ semántica; IA fuerte vs débil).

\T{Reformular como juego}
Jugadores · estado (todos) · turno (simultáneo $\to$ alternado) · acciones · transición (¿azar?) ·
terminal · utilidad. ¿Suma cero? ¿Info.\ perfecta?
\st{Minimax}: rival óptimo y hostil, peor caso, conservador. \st{Expectimax}: rival o entorno
probabilístico, maximiza lo esperado (racional si el rival no es malicioso). \st{Expectiminimax}: max + min + azar.
Árbol grande: profundidad límite + \textsc{Eval}, alfa-beta, profundización iterativa.
$n$ jugadores: tuplas, cada uno maximiza su componente; alianzas.

\T{Formular un problema de búsqueda}
Estado · inicial · $\textsc{Acciones}(s)$ · $\textsc{Resultado}(s,a)$ · meta · costo $c(s,a,s')$.
Tamaño $|S|$; $b$ = ramificación; $d$ = prof.\ solución; $m$ = prof.\ máx.; BFS genera $O(b^d)$.
Repetidos $\Rightarrow$ búsqueda en grafo (explorados); no pierde el óptimo: el costo futuro solo depende del estado.
Árbol vs grafo: un grafo finito con ciclos da un árbol infinito.

\T{Búsqueda no informada}
{\footnotesize
\begin{tabular}{@{}l@{\;\;}l@{\;\;}l@{\;\;}l@{\;\;}l@{}}
\toprule
 & Compl. & Ópt. & Tiempo & Espacio\\\midrule
BFS & sí & c.\,unif. & $b^d$ & $b^d$\\
UCS & sí ($\epsilon{>}0$) & sí & $b^{1+\lfloor C^*/\epsilon\rfloor}$ & ídem\\
DFS & no & no & $b^m$ & $bm$\\
IDS & sí & c.\,unif. & $b^d$ & $bd$\\
Bidir. & sí & sí & $b^{d/2}$ & $b^{d/2}$\\\bottomrule
\end{tabular}}
Frontera: BFS cola FIFO; UCS prioridad $g$; DFS pila LIFO; A* prioridad $g+h$.
\ojo{UCS/A*: prueba de meta al EXTRAER}, no al generar. IDS: $db+(d-1)b^2+\dots+b^d$
($b{=}10,d{=}5$: 123\,450 vs BFS 111\,110). BFS $\subset$ UCS (costos iguales); UCS = A* con $h{=}0$;
DFS = primero-el-mejor con $f{=}-$prof.

\T{A* y heurísticas}
$f(n)=g(n)+h(n)$. Voraz: $f=h$ (no óptima, en árbol cicla). $w$A*: $f=g+wh$, costo $\le wC^*$.
\begin{caja}
\st{Admisible}: $0\le h(n)\le h^*(n)$ (nunca sobreestima).\\
\st{Consistente}: $h(n)\le c(n,a,n')+h(n')$ $\forall$ arco, y $h(\text{meta})=0$.\\
Consistente $\Rightarrow$ admisible. Consistente $\Rightarrow$ $f$ no decrece $\Rightarrow$ al extraer, $g=g^*$
$\Rightarrow$ sin reexpansiones; A* en grafo óptimo.\\
Árbol + admisible $\Rightarrow$ óptimo: ancestro $n$ en frontera, $f(n)\le C^*<f(G_2)$.
\end{caja}
\st{Revisar}: $h^*$ hacia atrás desde $G$ (admisible: $h\le h^*$); consistencia \emph{arco por arco}.
\ojo{Grafo + inconsistente}: un nodo puede cerrarse con $g$ subóptimo (síntoma: $f$ baja). Arreglos:
reabrir nodos cerrados, búsqueda en árbol, pathmax $f(n')\gets\max(f(n'),f(n))$, $h$ consistente.
\st{Traza} (tabla): paso | extrae $(g,f)$ | frontera con $f$ | comentario (mejora, descarta, empate).
\st{Relajación} $\Rightarrow$ admisible y consistente. 8-puzzle: $h_1$ fuera de lugar (teletransporte),
$h_2$ Manhattan (mover aunque esté ocupada). Dominancia $h_2\ge h_1$: A* expande todo $f<C^*$ $\Rightarrow$
expande menos. $\max(h_1,h_2)$ sí; $h_1+h_2$ \emph{no} admisible. $b^*$: $N{+}1=\sum_{i=0}^{d}(b^*)^i$.
\st{Bidireccional}: parar si $\mu\le g^F_{\min}+g^B_{\min}$ (Dijkstra); MM: $pr=\max(f,2g)$, encuentro en $C^*/2$.

\T{CSP}
$(X,D,C)$; solución = completa y consistente. Restricciones: unaria, binaria, global (Alldiff), suave.
\st{Backtracking}: DFS, 1 variable por nivel ($d^n$ hojas, no $n!d^n$).
\st{FC}: al asignar $X$, borrar valores incompatibles de sus vecinos.
\st{Arco} $X\to Y$ consistente: $\forall x\in D_X\ \exists y\in D_Y$ compatible.
\st{AC-3}: si \textsc{Revise}$(X,Y)$ borra, encolar $(Z,X)$ $\forall Z\ne Y$; $O(cd^3)$.
Resultado: dominio vacío (sin sol.) / todos unitarios (sol.) / buscar.
\st{MRV}: menos valores (falla primero); \st{grado}: desempate; \st{LCV}: valor que menos restringe.
Árbol: $O(nd^2)$. Arco-consistente $\not\Rightarrow$ solución (triángulo, 2 colores).

\T{Búsqueda local}
Estados completos, memoria $O(1)$, incompleta. \st{Hill climbing}: máx.\ locales, mesetas, crestas;
reinicios: $1/p$ intentos. \st{Tabú}: mejor vecino \emph{no tabú} aunque empeore.
\st{Recocido} (maximizar): $\Delta E=f(n)-f(s)$; si $\Delta E>0$ acepta; si no, con prob.\ $e^{\Delta E/T}$;
$T\gets\alpha T$. $T$ alta = aleatorio; baja = hill climbing.
\st{GA}: aptitud, selección ($f_i/\sum f$ o torneo), cruza en un punto, mutación $P_m\approx1/L$;
$P(\text{sin mutar})=(1-P_m)^L$. No free lunch.

\T{Minimax y alfa-beta}
$V(s)=\textsc{Util}(s)$ si terminal; $\max$ en MAX; $\min$ en MIN. Tiempo $O(b^m)$, espacio $O(bm)$.
$\alpha$ = mejor valor de MAX en el camino; $\beta$ = mejor de MIN. Pasar $(\alpha,\beta)$ hacia abajo.
\ojo{En MIN: si $v\le\alpha$ podar. En MAX: si $v\ge\beta$ podar.} Actualizar $\alpha$ en MAX y $\beta$ en MIN.
Orden perfecto $O(b^{m/2})$ (hojas $b^{\lceil m/2\rceil}+b^{\lfloor m/2\rfloor}-1$); aleatorio $O(b^{3m/4})$.
La raíz es exacta; los nodos podados solo dan cotas.
\st{Evaluación}: $\textsc{Eval}(s)=\sum_i w_if_i(s)$; rápida, ordena terminales igual que la utilidad.
Gato: $3X_2+X_1-(3O_2+O_1)$. Efecto horizonte.

\T{Expectimax / expectiminimax}
\begin{caja}\small
$V(s)=\textsc{Util}(s)$ si es terminal (\textsc{Eval} en el límite)\\
MAX: $V(s)=\max_a\sum_{s'}P(s'|s,a)\,V(s')$\\
MIN: $V(s)=\min_a\sum_{s'}P(s'|s,a)\,V(s')$\\
Azar: $V(s)=\sum_i p_i\,V(s_i)$
\end{caja}
Costo $O(b^m n^m)$. Si $\rho=1$ es minimax. \ojo{Importa la escala}: solo invariante ante $ax+b$, $a>0$
(minimax: ante cualquier función creciente).
Desvío: $P(c'|c)=\rho[c'{=}c]+\frac{1-\rho}{|A|}$ (uniforme incl.\ $c$); \emph{declara el supuesto}.
Poda con utilidades en $[L,U]$: cota sup.\ $=\sum_{\text{vistos}}p_iv_i+(\sum_{\text{resto}}p_i)\,U$.

\T{Utilidad}
Lotería $L=[p,A;\,1{-}p,B]$; $A\succ B$ preferencia, $A\sim B$ indiferencia.
Axiomas: ordenabilidad, transitividad, continuidad, sustituibilidad, monotonicidad, descomponibilidad
$\Rightarrow\exists U$: $U(L)=\sum_ip_iU(S_i)$; elegir la de máxima utilidad esperada (MEU).
Allais: $B\succ A$ y $C\succ D$ $\Rightarrow$ $U(3k)>0.8U(4k)>U(3k)$: contradicción.
Equivalente cierto: $x$ con $U(x)=EU(L)$.
\end{multicols*}

\newpage
% =============================== REVERSO ===============================
{\large\bfseries\color{azul} IA · Examen 1 · Formulario (reverso): lógica}\par\vspace{2pt}
\begin{tcolorbox}[colback=azulclaro,colframe=azul,boxrule=0.5pt,arc=1pt,left=3pt,right=3pt,top=2pt,bottom=2pt,
  title={\bfseries SÍMBOLOS Y CÓMO SE LEEN \hfill TABLA DE VERDAD},fonttitle=\footnotesize,colbacktitle=azul,toptitle=0pt,bottomtitle=0pt]
\footnotesize
\begin{minipage}[t]{0.33\linewidth}
\begin{tabular}{@{}c@{\;\;}l@{}}
$\neg P$ & no $P$ (negación)\\
$P\wedge Q$ & $P$ y $Q$ (conjunción)\\
$P\vee Q$ & $P$ o $Q$ o ambos (inclusivo)\\
$P\imp Q$ & si $P$ entonces $Q$ (implicación)\\
$P\sii Q$ & $P$ si y solo si $Q$ (bicondicional)\\
$\top,\ \bot$ & \textit{True} (siempre V), \textit{False}\\
$\square$ & cláusula vacía = contradicción
\end{tabular}
\end{minipage}\hfill
\begin{minipage}[t]{0.33\linewidth}
\begin{tabular}{@{}c@{\;\;}l@{}}
$KB\models\alpha$ & vincula: $\alpha$ V en todo modelo de $KB$\\
$KB\vdash\alpha$ & deriva: se obtiene con reglas\\
$\alpha\equiv\beta$ & equivalentes: mismos modelos\\
$M(\alpha)$ & modelos donde $\alpha$ es V\\
$\forall x$,\ $\exists x$ & para todo $x$; existe algún $x$\\
$\{x/A\}$ & sustitución: $x$ se vuelve $A$\\
$\theta$ & unificador: $p\theta=q\theta$
\end{tabular}
\end{minipage}\hfill
\begin{minipage}[t]{0.31\linewidth}
\centering
\begin{tabular}{@{}cc|c@{\,}c@{\,}c@{\,}c@{\,}c@{}}
$P$ & $Q$ & $\neg P$ & $P{\wedge}Q$ & $P{\vee}Q$ & $P{\imp}Q$ & $P{\sii}Q$\\\hline
V & V & F & V & V & V & V\\
V & F & F & F & V & \textbf{F} & F\\
F & V & V & F & V & V & F\\
F & F & V & F & F & V & V
\end{tabular}\\[2pt]
\raggedright $P\imp Q$ solo es F con V$\imp$F. Precedencia: $\neg,\wedge,\vee,\imp,\sii$. $n$ símbolos: $2^n$ modelos.
\end{minipage}
\end{tcolorbox}
\begin{multicols*}{3}
\T{Definiciones semánticas}
\st{Modelo}: asignación V/F a cada símbolo. \st{Válida} (tautología): V en todo modelo.
\st{Satisfacible}: V en algún modelo. \st{Insatisfacible}: en ninguno. \st{Contingente}: en unos sí y en otros no.
\begin{caja}
$KB\models\alpha\iff M(KB)\subseteq M(\alpha)\iff KB\imp\alpha$ válida $\iff KB\wedge\neg\alpha$ insatisfacible.
\end{caja}
\st{Tell($f$)}: «ya lo sabía» ($KB\models f$); «no lo creo» ($KB\models\neg f$); «aprendí algo» (contingente).
\st{Ask($f$)}: sí / no / no sé. \st{Correcto}: $\vdash\subseteq\models$. \st{Completo}: $\models\subseteq\vdash$.
Monotonía: $KB\models\alpha\Rightarrow KB\wedge\beta\models\alpha$.

\T{Equivalencias clave}
{\small
$P\imp Q\equiv\neg P\vee Q\equiv\neg Q\imp\neg P$ (contrapositiva)\\
$P\sii Q\equiv(P\imp Q)\wedge(Q\imp P)$\\
$\neg(P\wedge Q)\equiv\neg P\vee\neg Q$;\ \ $\neg(P\vee Q)\equiv\neg P\wedge\neg Q$ (De Morgan)\\
$\neg\neg P\equiv P$;\ \ $\neg(P\imp Q)\equiv P\wedge\neg Q$\\
$P\vee(Q\wedge R)\equiv(P\vee Q)\wedge(P\vee R)$ (distribución)\\
$P\wedge(Q\vee R)\equiv(P\wedge Q)\vee(P\wedge R)$\\
$A\imp(B\imp C)\equiv(A\wedge B)\imp C$ (exportación)\\
$(A\wedge B)\imp C\equiv(A\imp C)\vee(B\imp C)$\\
$(A\imp C)\wedge(B\imp C)\equiv(A\vee B)\imp C$\\
$P\vee\neg P\equiv\top$;\ \ $P\wedge\neg P\equiv\bot$;\ \ $\bot\imp Q\equiv\top$}

\T{CNF proposicional}
1) quitar $\sii$; 2) quitar $\imp$; 3) meter $\neg$ (De Morgan, $\neg\neg$); 4) distribuir $\vee$ sobre $\wedge$.
Ej.: $B\sii(P\vee Q)\to(\neg B\vee P\vee Q)\wedge(\neg P\vee B)\wedge(\neg Q\vee B)$.
$(V\imp N)\imp R\to(V\vee R)\wedge(\neg N\vee R)$; $V\imp(N\imp R)\to\neg V\vee\neg N\vee R$.

\T{Inferencia}
\st{Modus ponens}: $\dfrac{P,\ P\imp Q}{Q}$ (correcto, no completo; completo con Horn).
\st{Resolución}: $\dfrac{\ell_1\vee\dots\vee p,\quad\neg p\vee m_1\vee\dots}{\ell_1\vee\dots\vee m_1\vee\dots}$
(correcta; completa por refutación).
\begin{caja}
\st{Refutación}: añadir $\neg\alpha$ a la KB, pasar todo a CNF, resolver hasta $\square$ $\Rightarrow KB\models\alpha$.
Si no sale $\square$: dar un \emph{modelo} de la KB con $\alpha$ falsa y decir qué hecho falta.
\end{caja}
\st{Horn}: $\le1$ literal positivo; \st{definida}: exactamente 1 ($p_1\wedge\dots\wedge p_k\imp q$); \st{meta}: 0.
Resolvente de Horn es Horn. $P\vee Q$ no es Horn.
\st{Encadenamiento adelante}: agenda + contador de premisas por regla; al llegar a 0, se infiere la
conclusión. Lineal. \st{Atrás}: de la meta a los hechos; evitar metas repetidas en la pila.
\st{DPLL}: terminación temprana, símbolo puro (un solo signo), cláusula unitaria, ramificar (V primero).
2-CNF: polinomial ($O(n^2)$ cláusulas posibles); 3-CNF: NP-completo.

\T{Lógica de primer orden}
Términos: constantes, variables, funciones $F(x)$. Átomos: predicados $P(t_1,..)$.
\begin{caja}
«Todo $P$ es $Q$»: $\forall x\,(P(x)\imp Q(x))$ \quad«Algún $P$ es $Q$»: $\exists x\,(P(x)\wedge Q(x))$\\
\ojo{Error típico}: $\forall$ con $\wedge$ (demasiado fuerte); $\exists$ con $\imp$ (demasiado débil).
\end{caja}
«Ningún $P$ es $Q$»: $\forall x\,(P(x)\imp\neg Q(x))$. «Solo uno»: $\exists x\,[P(x)\wedge\forall y\,(P(y)\imp y{=}x)]$.
$\neg\forall x\,P\equiv\exists x\,\neg P$; $\neg\exists x\,P\equiv\forall x\,\neg P$. $\exists y\forall x\models\forall x\exists y$ (no al revés).

\T{CNF en LPO (6 pasos)}
1) quitar $\sii,\imp$; 2) meter $\neg$ (cambia $\forall\leftrightarrow\exists$); 3) estandarizar variables;
4) \st{Skolem}: $\exists y$ dentro de $\forall x_1..x_k$ $\to$ $F(x_1..x_k)$; sin $\forall$ externo $\to$ constante;
5) quitar $\forall$; 6) distribuir $\vee$ sobre $\wedge$.
\ojo{$\exists$ en el antecedente de $\imp$ se vuelve $\forall$: sin Skolem.}
Ej.: $\forall x\,[\p{Persona}(x)\imp\exists y\,\p{Madre}(y,x)]\to\neg\p{Persona}(x)\vee\p{Madre}(M(x),x)$.

\T{Leer cláusulas en lenguaje natural}
$\neg A(x)\vee\neg B(x)\vee C(x)$ = «todo $x$ que es $A$ y $B$ es $C$». Literales negados = condiciones;
el positivo = conclusión; variables = «para todo». Hecho sin variables = afirmación concreta.

\T{Unificación}
Comparar argumento por argumento aplicando $\theta$ acumulado. Falla con constantes/funciones distintas
o por \st{ocurrencia} ($x$ vs.\ $F(x)$). Estandarizar aparte antes. MGU = el más general.
Ej.: $\p{Conoce}(\p{Juan},x)$, $\p{Conoce}(y,\p{Madre}(y))\to\{y/\p{Juan},x/\p{Madre}(\p{Juan})\}$.
$P(x,F(x))$, $P(F(y),y)$: falla (ocurrencia).

\T{Resolución en LPO: ejemplo tipo examen}
KB: (1) $\neg\p{Est}(x)\vee\neg\p{Toma}(x)\vee\p{Conoce}(x,a)$, (6) $\p{Est}(\p{bob})$, (7) $\p{Toma}(\p{bob})$.
(1)+(6) $\{x/\p{bob}\}$: $\neg\p{Toma}(\p{bob})\vee\p{Conoce}(\p{bob},a)$; +(7): $\p{Conoce}(\p{bob},a)$.
Refutación: añadir $\neg\alpha$, resolver hasta $\square$, anotar $\theta$ en cada paso; factorizar repetidos.
\st{Prolog}: Horn, SLD en profundidad, arriba$\to$abajo, izq.$\to$der.; recursión por la izquierda cicla.

\T{Ejemplos resueltos (repaso rápido)}
\st{Validez por tabla}: $(P\imp Q)\imp(\neg Q\imp\neg P)$ es válida (contrapositiva);
$(P\imp Q)\imp(Q\imp P)$ contingente (falla con $P{=}$F, $Q{=}$V); $P\wedge\neg P$ insatisfacible.
\st{Resolución prop.}: KB $\{A\imp(B\vee C),A,\neg B,\neg C\}$: $(\neg A\vee B\vee C)+A\to B\vee C$; $+\neg B\to C$; $+\neg C\to\square$.
\st{A* inconsistente (2025)}: $S{\to}A(1)$, $S{\to}B(4)$, $A{\to}C(2)$, $B{\to}C(1)$, $C{\to}G(3)$; $h(S,A,B,C,G)=5,3,2,1,0$.
Extrae S(0,5), A(1,4), C(3,4), B(4,6), G(6,6): costo 6 óptimo ($h^*(S)=6$). Viola $S\to A$: $5>1+3$.
\st{Expectiminimax}: azar $0.5\cdot\min(3,9)+0.5\cdot\min(6,4)=3.5$.
\st{Desvío}: atacar $2\rho-1$ vs bloquear $-(1-\rho)/2$: atacar si $\rho>1/3$.
\st{Recocido}: $\Delta E=-2$, $T=1$: $e^{-2}=0.135$; $T=10$: $0.82$.
\st{Unificación}: $\p{Ama}(x,\p{Padre}(x))$, $\p{Ama}(\p{Padre}(y),z)$: $\{x/\p{Padre}(y),\ z/\p{Padre}(\p{Padre}(y))\}$.
\st{Skolem}: $\exists x\forall y\,\p{Ama}(y,x)\to\p{Ama}(y,C)$; $\forall x[\forall y\,P(x,y)\imp\exists z\,Q(x,z)]\to\neg P(x,F(x))\vee Q(x,G(x))$.

\T{Recordatorios de examen}
Escribe la definición antes de usarla. Declara supuestos (desempates, azar). En «¿es óptima?», compara
con $h^*$ u otro camino. Justifica cada propiedad del ambiente. Revisa el tiempo: no te atores en una pregunta.
\end{multicols*}
\end{document}
